\section{参考文献}
引用别人的成果或其他公开的资料(包括网上查到的资料)必须按照规定的参考文献的表述方式在正文引用处和参考文献中均明确列出。正文引用处用方括号标示参考文献的编号，如[1][3]等；引用书籍还必须指出页码。参考文献按正文中的引用次序列出，其中

书籍的表述方式为：

[编号] 作者，书名，出版地：出版社，出版年.

参考文献中期刊杂志论文的表述方式为：

[编号] 作者，论文名，杂志名，卷期号：起止页码，出版年.

参考文献中网上资源的表述方式为：

[编号] 作者，资源标题，网址，访问时间（年月日）.
\begin{thebibliography}{99}
\bibitem{} 汤国生.基于SWOT分析法的大学生数学建模创新实践基地建设探索——以江苏科技大学为例[J].高教学刊,2023,9(11):53-56.DOI:10.19980/j.CN23-1593/G4.2023.11.013.
\bibitem{} 谢雨婧,韩惠丽.北师大版初中数学教材中数学建模的多维度分析[J].教学与管理,2023(09):73-76.
\bibitem{} 刘志梅.数学建模与高职数学教学的深度融合[J].佳木斯职业学院学报,2023,39(03):152-154.
\bibitem{} 许亚桃,吴立宝.基于Delphi-AHP高中数学建模教学评价指标体系的研究[J].内江师范学院学报,2023,38(02):113-119.DOI:10.13603/j.cnki.51-1621/z.2023.02.018.
\bibitem{} 杨本朝,石雅男,段乾恒,李光松,于刚.大学生数学建模竞赛开展全周期教学实践探究[J].大学教育,2023(04):44-46.
\bibitem{} 黄健,徐斌艳.国际视野下数学建模教与学研究的发展趋势——基于第14届国际数学教育大会的分析[J].数学教育学报,2023,32(01):93-98.
\begin{figure}[!h]
    \centering
    \includegraphics[width=.6\textwidth]{示例图1}
    \caption{示例图1}
    \label{fig:circuit-diagram}
\end{figure}

\begin{figure}[!h]
    \centering
    \includegraphics[width=.6\textwidth]{示例图2}
    \caption{示例图2}
    \label{fig:circuit-diagram}
\end{figure}
\end{thebibliography}